\section{Introduction}

\subsection{General concepts in SPC and ML}

\subsubsection{Statistical Process Control}

Statistical process control (SPC) is a collection of statistical techniques for monitoring process
stability and improving capability by reducing variability \citep{montgomery_statistical_2012}. Its
principal tool is the control chart, which plots a monitoring statistic against a control limit;
a value beyond the limit signals a possible out-of-control (OOC) departure from the in-control (IC)
state. Charts are judged by their run-lengths, usually through the average run-length (\ARL{}): a long IC \ARL{} ($\ARL_0$) means few false alarms and a short OOC \ARL{} ($\ARL_1$) fast detection \citep{ahmed_statistical_2017}. Memoryless charts use only the current observation, which makes them responsive to large shifts but insensitive to small ones; the Shewhart chart, for example, signals when the current sample mean falls outside three-sigma lower and upper control limits (\LCL{} and \UCL{}) around the IC mean \citep{ahmed_statistical_2017}. Memory-type charts such as the cumulative sum (CUSUM) \citep{page_continuous_1954} and the
exponentially weighted moving average (EWMA) \citep{roberts_control_1959} accumulate evidence and
suit small, persistent shifts; the multivariate EWMA (MEWMA) \citep{lowry_multivariate_1992} extends
this to correlated characteristics. These charts assume independent, normally distributed
observations, and when either assumption fails the attained $\ARL_0$ can differ substantially from its nominal value \citep{lowry_multivariate_1992,ahmed_statistical_2017}. These limitations motivate more flexible monitoring schemes, including nonparametric \citep{arslan_rankbased_2026} and machine learning (ML)-based control charts \citep{mukhtiar_enhanced_2025}.

\subsubsection{Machine Learning}

ML algorithms learn patterns from training data and so offer alternative ways to identify process shifts, faults and anomalies without a parametric process model; as labelled faults are rare, relevant methods model only normal behaviour
\citep{pimentel_review_2014}. An autoencoder is a neural network trained to reconstruct its input
through a low-dimensional bottleneck \citep{hinton_reducing_2006}. Once it is trained on IC data, large reconstruction errors signal departures such as anomalies or shifts. A convolutional autoencoder (CAE) extracts local patterns from short multivariate time-series windows
\citep{lecun_gradient-based_1998,honkela_stacked_2011}. The one-class support vector machine
(OC-SVM), an unsupervised extension of the support vector machine \citep{cortes_support-vector_1995},
encloses most normal data in a boundary and scores new observations against it \citep{scholkopf_estimating_2001}. On CAE features the two signals provide complementary information: the
score indicates whether a window's learned representation is typical of normal process behaviour, whereas the reconstruction error reflects how accurately the network reproduces the window.

\subsection{Problem statement and significance of the study}

Learned detectors can represent the nonlinear, serially dependent structure that weakens standard
control limits, and they are usually applied observation by observation against a
fixed threshold, without accumulating evidence or a run-length interpretation of their alarm rate
\citep{mabu_disaster_2019,biegel_deep_2022,guzman_sidewalk_2026}. This study places CAE
reconstruction information and OC-SVM anomaly-score information inside a memory-type MEWMA chart
and evaluates the result as a control chart: its limit is calibrated to a nominal $\ARL_0$ and its run-lengths are compared with those of a classical MEWMA calibrated to the same target, so any
benefit of learning is judged on the classical chart's own terms.

\subsection{Research aim and objectives}

The aim of this study is to develop and evaluate a hybrid ML-based MEWMA monitoring scheme that integrates a CAE and
an OC-SVM for detecting process shifts in nonlinear multivariate process data. The specific objectives of the study are to:
\begin{itemize}
\item construct a CAE and OC-SVM hybrid scheme that combines both signals in one monitoring vector;
\item investigate its run-length characteristics, with particular interest in small shifts;
\item evaluate the CAE and OC-SVM components using appropriate diagnostics;
\item compare the hybrid chart with a classical MEWMA under a common IC design target;
\item examine its performance across shift magnitudes and three smoothing constants; and
\item demonstrate the proposed hybrid framework on simulated data, including an industrial benchmark.
\end{itemize}

\subsection{Research outline}

The study is confined to unsupervised monitoring of simulated processes, so the true state is always
known: a nonlinear process built for the primary experiment and, as an external application, the
Extended Tennessee Eastman Process (TEP) benchmark. It covers the Phase~I regime, in which all components and limits are estimated from IC data, and the Phase~II regime, in which the frozen scheme is applied to independent sequences
\citep{ahmed_statistical_2017}. Section~2 reviews the literature on control charts and ML for process monitoring. Section~3 describes the experimental design, data generation, control charts, ML models and evaluation criteria. Section~4 presents the primary simulation results, Section~5 applies the approach to the TEP, and Section~6 summarises the findings, recommendations, limitations and future work.
